\documentclass[11pt,a4paper]{article}

% --- Pacotes Fundamentais ---
\usepackage[utf8]{inputenc}
\usepackage[T1]{fontenc}
\usepackage[brazil]{babel}
\usepackage{lmodern}
\usepackage[top=2.5cm, bottom=2.5cm, left=2.5cm, right=2.5cm]{geometry}
\usepackage{amsmath,amssymb}
\usepackage{graphicx}
\usepackage{booktabs}
\usepackage{listings}
\usepackage{xcolor}
\usepackage{fancyhdr}
\usepackage[most]{tcolorbox}
\usepackage{hyperref}

% --- Configurações de Cores e Links ---
\definecolor{primary}{RGB}{16, 54, 98}      % Azul institucional
\definecolor{secondary}{RGB}{41, 128, 185}  % Azul médio
\definecolor{boxbg}{RGB}{245, 248, 252}     % Fundo para caixas de instrução
\definecolor{boxframe}{RGB}{180, 205, 230}  % Borda das caixas de instrução
\definecolor{codebg}{RGB}{248, 249, 250}

\hypersetup{
    colorlinks=true,
    linkcolor=primary,
    citecolor=secondary,
    urlcolor=secondary
}

% --- Caixa de Instrução para os Alunos ---
\newtcolorbox{instrucao}{
    enhanced,
    colback=boxbg,
    colframe=boxframe,
    arc=3mm,
    boxrule=0.8pt,
    left=4mm, right=4mm, top=3mm, bottom=3mm,
    fonttitle=\bfseries\small\color{primary},
    title={Orientações para esta seção:},
    fontupper=\small\itshape\color{black!80}
}

% --- Configurações para Código-Fonte (Listings) ---
\lstset{
    backgroundcolor=\color{codebg},
    basicstyle=\ttfamily\footnotesize,
    breaklines=true,
    frame=single,
    rulecolor=\color{black!20},
    numbers=left,
    numberstyle=\tiny\color{black!40},
    keywordstyle=\color{primary}\bfseries,
    commentstyle=\color{green!50!black}\itshape,
    stringstyle=\color{orange!80!black}
}

% --- Cabeçalho e Rodapé ---
\pagestyle{fancy}
\fancyhf{}
\lhead{\footnotesize Projeto em Ciência e Inovação -- 2026/2}
\rhead{\footnotesize Relatório Técnico - Ciclo 3}
\cfoot{\thepage}

% ==============================================================================
% DADOS DO TRABALHO
% ==============================================================================
\title{
    \vspace{-1.0cm}
    \textbf{\Large [Título do Projeto / Solução]}\\[0.3cm]
    \large Relatório Técnico de Pesquisa e Desenvolvimento (P\&D)
}

\author{
    \textbf{Aluno 1} (Cartão: 000000) -- \texttt{aluno1@inf.ufrgs.br}\\
    \textbf{Aluno 2} (Cartão: 000000) -- \texttt{aluno2@inf.ufrgs.br}\\
    \textbf{Aluno 3} (Cartão: 000000) -- \texttt{aluno3@inf.ufrgs.br}\\[0.3cm]
    \normalsize Instituto de Informática -- Universidade Federal do Rio Grande do Sul (UFRGS)\\
    \normalsize Disciplina: Projeto em Ciência e Inovação (2026/2)\\
    \small Repositório do Projeto: \url{https://github.com/usuario/repositorio}
}

\date{Porto Alegre, \today}

% ==============================================================================
% INÍCIO DO DOCUMENTO
% ==============================================================================
\begin{document}

\maketitle
\thispagestyle{fancy}

\begin{abstract}
\noindent
\textbf{Resumo.} 
\begin{instrucao}
O resumo deve sintetizar o ciclo completo de P\&D (máximo de 250 palavras): (1) O problema real investigado; (2) As hipóteses centrais de solução; (3) A implementação computacional (estruturas de dados e algoritmos); e (4) A metodologia experimental, com os principais resultados e conclusões obtidas.
\end{instrucao}
[Insira aqui o texto do resumo executivo. Exemplo: Este relatório descreve o ciclo de P\&D conduzido para o problema de otimização de rotas de transporte coletivo a partir de dados abertos municipais. Formulou-se a hipótese de que a modelagem via grafo esparso em lista de adjacências combinada ao algoritmo de Dijkstra proporciona tempo sub-segundo para mais de 95\% das consultas reais. A solução foi implementada e validada através de baterias experimentais com 42 linhas urbanas, confirmando a hipótese com redução de 38\% no tempo de consulta frente à matriz de adjacência.]
\\[0.3cm]
\textbf{Palavras-chave:} Pesquisa e Desenvolvimento (P\&D); Método Científico; Hipóteses; Avaliação Experimental; Estruturas de Dados; Dados Abertos.
\end{abstract}

\vspace{0.3cm}
\hrule
\vspace{0.4cm}

% ------------------------------------------------------------------------------
\section{Contexto e Definição do Problema}

\begin{instrucao}
Apresente o cenário prático real que motivou a investigação e a formalização computacional do problema.
\end{instrucao}

Descreva o domínio de aplicação investigado (ex.: mobilidade urbana, saúde, alocação de recursos, transparência). Por que esse problema é relevante e demanda uma solução computacional? Quando pertinente, apresente a base de dados reais ou abertos utilizada: órgão provedor, volume de registros, principais atributos e formato.

\subsection{Formulação do Problema Computacional}
Como o cenário do mundo real foi formalizado como um problema computacional bem definido? Quais são as entradas (\textit{inputs}), restrições operacionais e saídas esperadas (\textit{outputs})?

% ------------------------------------------------------------------------------
\section{Formulação das Hipóteses de Solução}

\begin{instrucao}
Uma hipótese em computação experimental deve ser mensurável, precisa e passível de validação/falsificação. Defina o que o grupo esperava observar antes de realizar os experimentos.
\end{instrucao}

\subsection{Hipóteses Formuladas}
\begin{itemize}
    \item \textbf{Hipótese Principal ($H_1$):} Premissa central sobre a viabilidade, eficiência ou qualidade da solução proposta. \\
    \textit{Exemplo: ``A representação da malha por lista de adjacências associada ao algoritmo de Dijkstra permite consultas de menor caminho com tempo de resposta inferior a 50\,ms para o grafo da cidade.''}
    \item \textbf{Hipótese Secundária ou Comparativa ($H_2$):} Se aplicável, comparação entre abordagens ou estruturas alternativas. \\
    \textit{Exemplo: ``A ordenação prévia das demandas por uma heurística gulosa atinge custo de solução dentro de 10\% da solução exata, com consumo de memória 60\% menor.''}
\end{itemize}

\subsection{Métricas de Sucesso}
Quais métricas foram estabelecidas previamente para determinar se as hipóteses foram confirmadas ou refutadas? (Ex.: tempo de execução em milissegundos, consumo de memória, taxa de erro/acurácia, custo de processamento).

% ------------------------------------------------------------------------------
\section{Desenvolvimento e Implementação da Solução}

\begin{instrucao}
Descreva a realização técnica: como as entidades reais foram mapeadas em estruturas de dados e quais algoritmos foram desenvolvidos.
\end{instrucao}

\subsection{Modelagem Computacional e Arquitetura}
Explique a representação formal em estruturas de dados:
\begin{itemize}
    \item Vértices $V$ e arestas $E$ no caso de grafos (o que representam no domínio real?);
    \item Estruturas de indexação rápida (tabelas hash, árvores);
    \item Filas de prioridade/heaps para seleção de menor custo.
    \item entre outros.
\end{itemize}

\subsection{Algoritmos e Análise de Custo/Complexidade}
Apresente a lógica dos algoritmos implementados para testar as hipóteses. Justifique as escolhas e apresente uma análise de custo/complexidade da solução.

\subsection{Tratamento e Pré-processamento dos Dados}
Quanod pertinente, descreva o pipeline de obtenção e preparação dos dados brutos: limpeza de registros nulos, normalização, filtros de consistência e geração das estruturas necessárias para os experimentos.

% ------------------------------------------------------------------------------
\section{Projeto Experimental e Metodologia de Validação}

\begin{instrucao}
Descreva com rigor o protocolo experimental para assegurar a reprodutibilidade dos resultados.
\end{instrucao}

\subsection{Ambiente Experimental e Variáveis}
\begin{itemize}
    \item \textbf{Ambiente:} Especificações de hardware (processador, memória RAM) e software (sistema operacional, versão da linguagem).
    \item \textbf{Variáveis Independentes:} O que foi manipulado nos experimentos (tamanho da entrada $N$, densidade do grafo, etc.).
    \item \textbf{Variáveis Dependentes:} O que foi medido como resposta (tempo de execução, memória de pico, etc.).
    \item \textbf{Variáveis de Controle:} Condições mantidas constantes (número de repetições, descarte de \textit{warm-up}, condições de teste).
\end{itemize}

\subsection{Cenários e Protocolo de Coleta}
Descreva os cenários de teste aplicados (dados sintéticos controlados vs. dados reais completos) e, quando pertinente, a metodologia estatística de consolidação (cálculo de médias, desvios-padrão).

% ------------------------------------------------------------------------------
\section{Resultados Experimentais e Discussão}

\begin{instrucao}
Apresente os dados quantitativos em tabelas e gráficos, confrontando-os diretamente com as hipóteses formuladas na Seção 2.
\end{instrucao}

\subsection{Apresentação dos Dados Coletados}
Apresente os resultados por meio de tabelas estruturadas e gráficos de desempenho.

\begin{table}[ht]
\centering
\caption{Desempenho comparativo de tempo e memória em função do tamanho da entrada $N$}
\label{tab:experimentos}
\resizebox{\textwidth}{!}{%
\begin{tabular}{@{}ccccc@{}}
\toprule
\textbf{Entrada ($N$)} & \textbf{Tempo } & \textbf{Tempo} & \textbf{Desvio} &
\textbf{Pico de} \\
    & \textbf{Abordagem 1 (ms)} & \textbf{Abordagem 2 (ms)} & \textbf{Padrão} &
\textbf{Memória} \\ \midrule
1.000                  & 8,2                             & 4,1                             & $\pm$ 0,3 ms           & 12 MB                    \\
5.000                  & 45,7                            & 18,3                            & $\pm$ 1,1 ms           & 24 MB                    \\
10.000                 & 110,4                           & 39,6                            & $\pm$ 2,4 ms           & 48 MB                    \\ \bottomrule
\end{tabular}%
}
\end{table}

\subsection{Validação ou Refutação das Hipóteses}
\begin{itemize}
    \item \textbf{Hipótese $H_1$:} Foi confirmada ou refutada? Indique explicitamente as evidências empíricas obtidas.
    \item \textbf{Complexidade Empírica vs. Teórica:} A curva de crescimento do tempo nos gráficos acompanhou a ordem $\mathcal{O}$ calculada?
\end{itemize}

\subsection{Discussão Técnica e Limitações Observadas}
Quais gargalos de desempenho e limitações de implementação surgiram ao longo dos experimentos? Ocorreram comportamentos inesperados?

% ------------------------------------------------------------------------------
\section{Análise de Impacto e Transferência Tecnológica}

\begin{instrucao}
Esta seção avalia o valor gerado além do laboratório. Como a solução computacional impacta a sociedade, a gestão de dados e qual é o caminho para sua aplicação prática?
\end{instrucao}

\subsection{Análise Multidimensional de Impacto}
\begin{itemize}
    \item \textbf{Impacto Social e Cidadania:} De que forma a solução melhora o serviço prestado, a qualidade de vida da comunidade ou a transparência pública? Quem são os beneficiários diretos e indiretos?
    \item \textbf{Impacto Econômico e Eficiência Operacional:} Quais ganhos de produtividade, otimização de tempo ou redução de custos a automação/algoritmo proporciona em comparação aos processos manuais ou legados?
    \item \textbf{Dimensão Ética, Legal e Governança de Dados (LGPD):}
    \begin{itemize}
        \item Há uso de dados pessoais ou sensíveis? Como a privacidade e a anonimização foram tratadas em conformidade com a LGPD?
        \item Há riscos de viés, discriminação algorítmica ou exclusão decorrentes da forma como os dados foram coletados ou ponderados?
        \item A solução garante rastreabilidade e explicabilidade das decisões/saídas geradas?
    \end{itemize}
    \item \textbf{Impacto Ambiental:} Há contribuição direta ou indireta para sustentabilidade (ex.: redução de emissões por otimização logística, uso eficiente de energia e recursos computacionais)?
\end{itemize}

\subsection{Viabilidade de Implantação e Transferência Tecnológica}
\begin{itemize}
    \item \textbf{Aderência ao Setor de Aplicação:} A solução possui potencial de ser adotada por órgãos públicos (prefeituras, secretarias municipais), empresas ou organizações da sociedade civil? Quem seriam os parceiros naturais para transferência de tecnologia?
    \item \textbf{Desafios para Ambiente de Produção:} Quais são os requisitos mínimos de infraestrutura (servidores, bancos de dados, conectividade) e os gargalos para migrar o protótipo experimental para um sistema contínuo em operação?
    \item \textbf{Manutenibilidade e Sustentabilidade da Base:} Como os dados devem ser atualizados e higienizados periodicamente para garantir que a solução permaneça útil ao longo do tempo?
\end{itemize}

% ------------------------------------------------------------------------------
\section{Conclusões e Lições Aprendidas do Processo de P\&D}

\begin{instrucao}
Sintetize os aprendizados do ciclo de P\&D: o que deu certo, o que não funcionou e quais seriam as próximas investigações.
\end{instrucao}

\subsection{Síntese dos Resultados Técnicos}
O que os resultados comprovaram a respeito da viabilidade e eficiência da solução computacional desenvolvida para o problema prático?

\subsection{O Que Funcionou vs. O Que Não Funcionou}
Relate com transparência as estratégias bem-sucedidas e as tentativas iniciais que falharam ou precisaram ser reformuladas ao longo do projeto.

\subsection{Próximos Passos de Pesquisa e Desenvolvimento}
Quais melhorias algorítmicas, heurísticas mais refinadas ou novas estruturas de dados deveriam ser exploradas em uma próxima iteração de P\&D?

% ------------------------------------------------------------------------------
\begin{thebibliography}{99}

\bibitem{dadosabertos}
[Nome da Base / Portal de Dados Abertos]. 
\textit{[Título do Conjunto de Dados Utilizado]}. Disponível em: \url{https://dados.gov.br}. Acesso em: \today.

\bibitem{cormen}
T. H. Cormen, C. E. Leiserson, R. L. Rivest, and C. Stein. 
\textit{Introduction to Algorithms}. 3rd ed. MIT Press, 2009.

\bibitem{metodologia}
[Referência Metodológica de Computação Experimental / Métodos Científicos].

\end{thebibliography}

% ------------------------------------------------------------------------------
\appendix
\section{Guia de Reprodutibilidade}

\begin{instrucao}
Forneça os comandos e requisitos necessários para executar o código e reproduzir os experimentos descritos.
\end{instrucao}

\begin{lstlisting}[language=bash, caption={Instruções de Execução dos Experimentos}]
# 1. Clonar o repositorio
git clone https://github.com/usuario/repositorio.git
cd repositorio

# 2. Instalar dependencias
pip install -r requirements.txt

# 3. Executar preparacao dos dados e baterias de teste
python main.py --preparar-dados
python main.py --executar-experimentos --gerar-graficos
\end{lstlisting}

\end{document}
